\section{Perception under causal observability}
\label{sec:perception}

\Cref{sec:causal-obs} of the framework introduced the perception
side of the architecture: the observation map $\obsmap$ of
\cref{def:obsmap}, the cognitive boundary $D(\jointvec)$ of
\cref{def:cognitive-boundary}, and the causal-observability
condition of \cref{def:causal-obs} that adds an explicit deadline
to the partial-observation reading. The present section specializes
that abstract framework to a concrete manipulation pipeline and
develops it in four steps: \cref{sec:perception-wrist-fov} gives
the explicit FOV-cone form of the cognitive boundary for a
wrist-mounted camera; \cref{sec:obs-action-coupling} formalizes the
observation--action coupling that is implicit in any wrist-mounted
sensor and identifies its place in the bi-level scaffold of
\cref{sec:bilevel}; \cref{sec:perception-timeliness} sharpens the
latency-budget inequality \eqref{eq:timeliness} into a per-phase
condition aligned with the five-phase decomposition of
\cref{sec:phase-cycle}; and \cref{sec:perception-three-modalities}
specializes the single observation map to three sensory modalities
(vision, proprioception, haptics), each with its own cognitive
boundary, latency profile, and best-served phases of the cycle. The
stochastic extension in which $D(\jointvec)$ is replaced by a
probability-of-detection field is deferred to
\cref{app:stochastic-fov}.

% ---------------------------------------------------------------
\subsection{Cognitive boundary of a wrist-mounted sensor}
\label{sec:perception-wrist-fov}

\Cref{def:cognitive-boundary} of \cref{sec:causal-obs} introduced
the cognitive boundary $D(\jointvec) \subseteq \mathbb{R}^{3}
\times SE(3)$ of the observation map $\obsmap$ as the
configuration-dependent subset of object positions and
orientations distinguishable from the sensor reading. The
abstract definition leaves the geometry of $D(\jointvec)$
unspecified. The present subsection specializes
$D(\jointvec)$ to the wrist-mounted-camera case used throughout
\cref{sec:perception}, with the camera frame rigidly attached to
the manipulator's end-effector and the field of view modeled as a
right circular cone with apex at the camera optical center.

\paragraph{Wrist-mounted FOV cone.}
Let $T_{\mathrm{cam}}(\jointvec) \in SE(3)$ denote the
camera-to-world transform at joint configuration $\jointvec$,
obtained by composing the manipulator forward kinematics from
the base to the end-effector with the fixed end-effector-to-camera
mounting transform. The camera FOV is parametrized by the
half-angle $\theta \in (0, \pi/2)$ of the cone about the optical
axis $\hat{z}_{\mathrm{cam}}$ and by the near and far clipping
planes $z_{\mathrm{near}}, z_{\mathrm{far}} \in \mathbb{R}_{>0}$
with $z_{\mathrm{near}} < z_{\mathrm{far}}$. The FOV cone in the
camera frame is the truncated cone
\begin{equation}
\label{eq:fov-cone-cam}
  C_{\mathrm{cam}}
  \;=\;
  \bigl\{\, p^{\mathrm{cam}} \in \mathbb{R}^{3}
  \;:\;
  z_{\mathrm{near}} \le p^{\mathrm{cam}}_{z} \le z_{\mathrm{far}},
  \;\;
  \|p^{\mathrm{cam}}_{xy}\| \le p^{\mathrm{cam}}_{z} \tan\theta
  \,\bigr\}.
\end{equation}

\begin{definition}[Wrist-mounted FOV cognitive boundary]
\label{def:wrist-fov}
For a wrist-mounted camera with mounting transform
$T_{\mathrm{cam}}(\jointvec)$, half-angle $\theta$, and clipping
planes $(z_{\mathrm{near}}, z_{\mathrm{far}})$, the cognitive
boundary $D(\jointvec) \subseteq \mathbb{R}^{3} \times SE(3)$ of
\cref{def:cognitive-boundary} is the set
\begin{equation}
\label{eq:wrist-fov-D}
  D(\jointvec)
  \;=\;
  \bigl\{\, (x, R) \in \mathbb{R}^{3} \times SE(3)
  \;:\;
  T_{\mathrm{cam}}(\jointvec)^{-1}\, x \;\in\; C_{\mathrm{cam}}
  \,\bigr\},
\end{equation}
i.e., the set of object poses whose position $x$ lies inside the
camera-frame truncated cone $C_{\mathrm{cam}}$ of
\eqref{eq:fov-cone-cam} after pull-back through
$T_{\mathrm{cam}}(\jointvec)$. The orientation component $R$ is
unconstrained at the level of the FOV containment test; finer
recoverability of $R$ from a single image is treated by the
modality-specific noise model in \cref{sec:perception}.
\end{definition}

\paragraph{Occlusion shrinks $D(\jointvec)$.}
The set $D(\jointvec)$ in \eqref{eq:wrist-fov-D} accounts only
for the FOV containment test. An object inside the cone but
whose line of sight from the camera optical center is blocked by
other geometry in the scene does not project onto the image
plane and is therefore not distinguishable from the sensor
reading. The deterministic FOV cone of \cref{def:wrist-fov} thus
overestimates the true cognitive boundary, and occlusion further
shrinks it. The exact treatment, in which $D(\jointvec)$ is
replaced by a probability-of-detection field over
$\mathbb{R}^{3} \times SE(3)$ that integrates FOV containment,
occlusion, and sensor noise, is deferred to
\cref{app:stochastic-fov}.

\paragraph{Bi-manual dependence on both configurations.}
For a Stage IV-type bi-manual scene in which a hammer, nail, and
block are manipulated by a primary arm in the presence of a
second arm, the second arm's links can occlude the workpiece
from the primary arm's wrist camera. The FOV cone is determined
by the primary arm's configuration $\jointvec_{\mathrm{self}}$
through $T_{\mathrm{cam}}(\jointvec_{\mathrm{self}})$, but the
occluding geometry depends on the second arm's configuration
$\jointvec_{\mathrm{other}}$. The cognitive boundary is
therefore a function of both,
$D(\jointvec_{\mathrm{self}}, \jointvec_{\mathrm{other}})$,
and the architectural plan-feasibility check of
\cref{def:causal-obs} will be evaluated against the joint
configuration of both arms. The detailed probabilistic model of
the joint dependence is deferred to \cref{app:stochastic-fov}.

\begin{example}[Wrist FOV for a parallel-jaw end-effector]
\label{ex:wrist-fov-pj}
Consider a parallel-jaw end-effector with a wrist camera mounted
just above the gripper such that the optical axis
$\hat{z}_{\mathrm{cam}}$ points along the approach direction of
the jaws. Let the mounting transform at a typical pre-grasp
configuration $\jointvec_{0}$ place the camera optical center at
height $h$ above the table, with $\hat{z}_{\mathrm{cam}}$
pointing vertically downward. The world-frame footprint of the
FOV cone on the table is then the annular disk obtained by
intersecting $C_{\mathrm{cam}}$ pulled forward through
$T_{\mathrm{cam}}(\jointvec_{0})$ with the table plane: a disk
of radius $h \tan\theta$ centered under the camera, restricted to
$z_{\mathrm{near}} \le h \le z_{\mathrm{far}}$. As the arm swings
in the null-space direction of \cref{sec:hierarchical-tn} during
phase 1, $T_{\mathrm{cam}}(\jointvec)$ rotates and the world-frame
footprint sweeps over the workspace; the cognitive boundary
$D(\jointvec)$ of \cref{def:wrist-fov} sweeps along with it.
\end{example}

% ---------------------------------------------------------------
\subsection{Observation--action coupling}
\label{sec:obs-action-coupling}

The causal-observability framework of \cref{sec:causal-obs}
exposed a structural feature of manipulation perception that
distinguishes it from passive partial-observation settings: the
wrist-mounted sensor rides the end-effector, so the cognitive
boundary $D(\jointvec)$ of \cref{def:cognitive-boundary} is
itself a function of the control history. The
observation--action paragraph at the end of
\cref{sec:causal-obs} stated this informally; the present
subsection formalizes it as the \emph{observability lift} of
\cref{def:obs-action-coupling} below and records the
architectural consequence in
\cref{prop:active-perception}. The order parameter
$\rcontact$ is not modified by this subsection: the contact-side
interface remains the only named scalar of the framework, and
the observability lift acts on the perception-side map
$\obsmap$ of \cref{def:obsmap} through the configuration
$\jointvec(t)$ that the controller realizes.

\begin{definition}[Observability lift]\label{def:obs-action-coupling}
For a control history $u_{0:t-1}$ producing the trajectory
$\jointvec(\cdot)$ via the joint-space dynamics of
\cref{sec:bilevel}, the \emph{observability lift} is the
composite map
\[
  u_{0:t-1}
  \;\longmapsto\;
  \jointvec(t)
  \;\longmapsto\;
  D\!\bigl(\jointvec(t)\bigr),
\]
where the first arrow integrates the dynamics under
$u_{0:t-1}$ and the second arrow evaluates the cognitive
boundary of \cref{def:cognitive-boundary} at the realized
configuration. The image of the lift is the cognitive boundary
realized at time $t$ along the trajectory induced by the
control history.
\end{definition}

\begin{proposition}[Active perception in the bi-level scaffold]
\label{prop:active-perception}
Under the bi-level scaffold of \cref{def:bilevel}, the outer
level may include in the least-action term $u_{a}$ of
\cref{def:three-term} a perception-augmenting contribution
whose role is to enlarge $D(\jointvec(t))$ at upcoming
phase transitions identified by the operator algebra
$\phaseop{0}, \ldots, \phaseop{4}$ of
\cref{sec:operator-algebra}, subject to the timeliness
condition \eqref{eq:timeliness}.
\end{proposition}

\begin{proof}
By direct substitution: $u_{a}$ enters the control input
through \cref{def:three-term} and acts on $\jointvec(t)$ via
the dynamics of \cref{sec:bilevel}; the cognitive boundary
$D(\jointvec(t))$ depends on $\jointvec(t)$ by
\cref{def:cognitive-boundary}; therefore an outer-level choice
of $u_{a}$ that steers $\jointvec(t)$ toward configurations
with larger $D(\jointvec(t))$ enlarges the realized image of
the observability lift of
\cref{def:obs-action-coupling} at the upcoming phase boundary.
\end{proof}

\paragraph{Relation to active perception.}
\Cref{prop:active-perception} is the manipulation specialization
of active perception in the sense of \citet{Bajcsy1988}: the
sensor is rigidly attached to the controlled kinematic chain,
so every action moves the sensor, and active perception is
automatic rather than an additional control mode layered on a
fixed sensor pose. The same observability lift is in principle
available to a free-standing observer that pans a camera, but
in the manipulation setting the lift is realized by the same
$u_{a}$ that the bi-level scaffold of \cref{def:bilevel}
already issues for task purposes; perception augmentation and
task control share an actuation channel.

\paragraph{Forward references.}
\Cref{sec:perception-three-modalities} catalogs which sensory
modalities benefit from the observability lift of
\cref{def:obs-action-coupling}: vision benefits strongly,
because the FOV cone of a wrist-mounted camera is reoriented
by every motion of the arm; haptics benefits moderately, since
contact-patch sensing is enlarged by reorienting the fingertip
across the contact face; proprioception does not benefit, since
the joint-encoder map is independent of $\jointvec$.

% ---------------------------------------------------------------
\subsection{Information timeliness}
\label{sec:perception-timeliness}

\Cref{def:causal-obs} of \cref{sec:causal-obs} states the latency
budget \eqref{eq:timeliness},
$\tau_{\mathrm{sensor}} + \tau_{\mathrm{compute}} <
\tau_{\mathrm{deadline}}$, as a single global condition. The
five-phase decomposition \cref{def:five-phase} of
\cref{sec:phase-cycle} forces a finer reading: each phase
$k \in \{1, \ldots, 5\}$ has its own controlling event---either a
$\swsurf$-crossing in the sense of \cref{lem:two-crossing} or a
phase-internal waypoint---which sets a phase-local deadline
$\tau_{\mathrm{deadline}}^{(k)}$. The latency-budget inequality
\eqref{eq:timeliness} must hold separately at every phase, since
a single overrun at any phase invalidates the cycle through the
operator composition of \cref{thm:joint-loop}.

\paragraph{Per-phase deadlines.}
\Cref{tab:per-phase-deadlines} collects the controlling event and
the qualitative scale of $\tau_{\mathrm{deadline}}^{(k)}$ phase by
phase. The bang-arc phases 2 and 4, at which the spectral-kick
torque $u_{s}$ of \cref{def:three-term} drives $\specgap$ across
the threshold $\varepsilon$, set the tightest deadlines: the
estimate must reach the controller within the window in which the
$\swsurf$-crossing remains actionable. The out-of-contact reaches
1 and 5 set looser deadlines determined by the next pose
waypoint; the in-contact transport phase 3 sets an intermediate
deadline determined by the next force-closure correction along
the singular arc.

\begin{table}[h]
\centering
\small
\begin{tabular}{@{}cllc@{}}
\toprule
$k$ & Phase & Controlling event & Scale of $\tau_{\mathrm{deadline}}^{(k)}$ \\
\midrule
1 & pick (approach)            & next pose waypoint along $\phaseop{0}$        & loose \\
2 & load (contact)             & $\swsurf$-crossing at $t_{1}$                 & very tight \\
3 & use (transport)            & next force-closure correction on $(t_{1}, t_{4})$ & intermediate \\
4 & unload (release)           & $\swsurf$-crossing at $t_{4}$                 & very tight \\
5 & place (retreat)            & rest configuration at $t = T$                 & loose \\
\bottomrule
\end{tabular}
\caption{Per-phase deadlines for the latency-budget inequality
\eqref{eq:timeliness}. Phases 2 and 4, at which the bang torque
$u_{s}$ of \cref{def:three-term} must arrive at the right
instant, set the tightest deadlines; phases 1 and 5 are out-of-contact
reaches with the loosest deadlines; phase 3 sets an intermediate
deadline along the singular transport arc.}
\label{tab:per-phase-deadlines}
\end{table}

\begin{proposition}[Per-phase timeliness]
\label{prop:per-phase-timeliness}
The manipulation pipeline is causally observable, in the sense of
\cref{def:causal-obs}, on a full pick--place cycle if and only if
the latency-budget inequality \eqref{eq:timeliness} holds with
$\tau_{\mathrm{deadline}}$ replaced by
$\tau_{\mathrm{deadline}}^{(k)}$ for every phase
$k \in \{1, \ldots, 5\}$ of \cref{def:five-phase}.
\end{proposition}

\begin{proof}
Each phase has a phase-determined event whose deadline
$\tau_{\mathrm{deadline}}^{(k)}$ enters \eqref{eq:timeliness}
locally; failing the inequality at some phase $k$ produces a
state estimate that arrives after the corresponding event, hence
fails \cref{def:causal-obs}~(ii) on that phase. By the joint-space
loop theorem \cref{thm:joint-loop}, the cycle composes through
the operators $\phaseop{0}, \ldots, \phaseop{4}$, so a phase-local
failure cascades to the cycle. Conversely, holding the inequality
at every phase yields a cycle-wide causally observable trajectory.
\end{proof}

\paragraph{Parallel with the spectral-gap margin.}
The strict inequality \eqref{eq:timeliness} is the temporal
analog of the spectral-gap margin condition $\specgap \ge
\varepsilon$ of \cref{sec:specgap}, in the precise sense already
noted in \cref{sec:causal-obs}: both are quantitative margins the
controller maintains, the spectral-gap margin during in-contact
transport (phase 3) and the latency margin during perception at
every phase. Both degrade gracefully---a small shrinkage of the
margin does not break feasibility---and both are state-dependent
quantities the controller actively manages rather than
once-and-for-all design parameters.

\paragraph{Forward references.}
\Cref{sec:perception} treats three sensory modalities (vision,
proprioception, haptics), each with a distinct
$(\tau_{\mathrm{sensor}}, \tau_{\mathrm{compute}})$ profile.
The per-phase deadline schedule of
\cref{tab:per-phase-deadlines} determines, for each modality,
which phases its profile is compatible with: a modality with a
large $\tau_{\mathrm{sensor}} + \tau_{\mathrm{compute}}$ may
serve phases 1 and 5 but not the bang-arc phases 2 and 4, while a
modality with a small latency profile may serve every phase
including the $\swsurf$-crossings.

% ---------------------------------------------------------------
\subsection{Three sensory modalities}
\label{sec:perception-three-modalities}

The observation map $\obsmap$ of \cref{def:obsmap} in
\cref{sec:causal-obs} was introduced as a single deterministic
mapping from manipulator state to sensor reading. A real
manipulation pipeline operates on multiple sensors in parallel,
each with its own cognitive boundary in the sense of
\cref{def:cognitive-boundary} and its own latency profile in the
sense of \cref{eq:timeliness}. The present subsection
specializes $\obsmap$ to three modalities: a wrist-mounted
RGB(-D) camera, joint encoders together with joint-torque
sensors, and a fingertip tactile sensor. Each modality is
described by its own observation map---written
$\obsmap_{v}$ (vision), $\obsmap_{p}$ (proprioception), and
$\obsmap_{h}$ (haptics)---with its own per-configuration
cognitive boundary
$D_{v}(\jointvec)$, $D_{p}(\jointvec)$, $D_{h}(\jointvec)$ and
its own qualitative position on the latency budget of
\eqref{eq:timeliness}. The deliverable of the subsection is
\cref{tab:modality-comparison}, which catalogs the three
modalities along the dimensions of cognitive boundary, latency
profile, best-served phases of the five-phase cycle of
\cref{def:five-phase}, and coupling to the active-perception
mechanism of \cref{sec:obs-action-coupling}.

\begin{table}[h]
\centering
\small
\begin{tabular}{@{}p{0.13\linewidth}p{0.21\linewidth}p{0.13\linewidth}p{0.16\linewidth}p{0.21\linewidth}@{}}
\toprule
Modality & Cognitive boundary $D(\jointvec)$
         & Latency profile (qualitative)
         & Best-served phases
         & Active-perception coupling \\
\midrule
Vision ($\obsmap_{v}$)
  & wrist-camera FOV cone of \cref{def:wrist-fov}
  & moderate (image read-out plus neural decoder)
  & 1, 3, 5
  & strong (camera rides end-effector) \\
Proprioception ($\obsmap_{p}$)
  & all of joint space (no FOV restriction)
  & low (kHz-rate encoder bus)
  & 2, 4
  & none (encoder map independent of $\jointvec$) \\
Haptics ($\obsmap_{h}$)
  & immediate fingertip contact patch (empty out-of-contact)
  & low (analog contact transduction)
  & 2, 3, 4
  & moderate (fingertip must touch the object) \\
\bottomrule
\end{tabular}
\caption{Specialization of the observation map $\obsmap$ of
\cref{def:obsmap} to three sensory modalities. Cognitive
boundaries $D(\jointvec)$ are taken in the sense of
\cref{def:cognitive-boundary}; latency profiles are positioned on
the budget \eqref{eq:timeliness} of \cref{sec:causal-obs}
qualitatively; best-served phases refer to the five-phase
cycle of \cref{def:five-phase}; active-perception coupling refers
to the observability lift of \cref{def:obs-action-coupling} in
\cref{sec:obs-action-coupling}.}
\label{tab:modality-comparison}
\end{table}

\paragraph{Vision.}
The vision modality is realized by a wrist-mounted RGB or RGB-D
camera whose mounting transform $T_{\mathrm{cam}}(\jointvec)$
rides the end-effector. The corresponding observation map
$\obsmap_{v}$ takes a scene state to the image (and, for RGB-D,
depth) read off the sensor at configuration $\jointvec$, and the
cognitive boundary $D_{v}(\jointvec)$ is the wrist-camera FOV
cone of \cref{def:wrist-fov}, with occlusion further shrinking
$D_{v}$ as discussed there and treated probabilistically in
\cref{app:stochastic-fov}. The latency profile is moderate:
$\tau_{\mathrm{sensor}}$ includes image read-out and any
demosaicing or rectification, and $\tau_{\mathrm{compute}}$
includes the neural pose-decoder forward pass that maps an image
to an object-pose estimate. Vision is best-served in the
out-of-contact reaches (phases 1 and 5) and the in-contact
transport (phase 3), where the controller must track an object
or end-effector pose against a $\swsurf$-distant deadline. The
modality is strongly coupled to active perception in the sense of
\cref{prop:active-perception}: every motion of the arm reorients
$T_{\mathrm{cam}}(\jointvec)$ and therefore $D_{v}(\jointvec)$,
so the outer level of the bi-level scaffold may steer the
configuration toward larger $D_{v}(\jointvec)$ at upcoming phase
boundaries identified by the operator algebra
$\phaseop{0}, \ldots, \phaseop{4}$.

\paragraph{Proprioception.}
The proprioception modality is realized by joint encoders---and,
on manipulators so equipped, joint-torque sensors at each
actuator. The corresponding observation map $\obsmap_{p}$ takes
the manipulator state to the encoder reading $\jointvec$
(together with $\dot{\jointvec}$ from differentiation or a
dedicated tachometer channel, and joint torques where present).
The cognitive boundary $D_{p}(\jointvec)$ is essentially all of
joint space: encoders observe $\jointvec$ itself, with no FOV
restriction and no occlusion. The latency profile is the lowest
of the three modalities, with the encoder bus running at
several orders of magnitude above the slowest closed-loop control
rate, so $\tau_{\mathrm{sensor}} + \tau_{\mathrm{compute}}$ is
small relative to any phase-boundary deadline. Proprioception is
best-served in the bang arcs at the $\swsurf$-crossings (phases 2
and 4) of \cref{lem:two-crossing}, where fast joint-velocity
feedback is the prerequisite for detecting that the order
parameter $\rcontact$ has crossed the threshold of
\cref{sec:Sigma}. Proprioception is not coupled to active
perception: $D_{p}(\jointvec) = D_{p}$ is independent of the
configuration, and the observability lift of
\cref{def:obs-action-coupling} acts trivially on $\obsmap_{p}$.

\paragraph{Haptics.}
The haptics modality is realized by a fingertip tactile sensor or
a fingertip force sensor, sensing pressure or shear distribution
across the contact patch. The corresponding observation map
$\obsmap_{h}$ takes the manipulator-and-object state to the local
contact-patch reading at the fingertip. The cognitive boundary
$D_{h}(\jointvec)$ is the immediate contact patch at
$\jointvec$: empty when the fingertip is out of contact (so that
the modality reports nothing distinguishable), non-empty during
the in-contact phases identified by $\rcontact \ge 1$ in the
sense of \cref{sec:causal-obs}. The latency profile is low: the
contact transduction stage is analog and short, and the
$\tau_{\mathrm{compute}}$ stage is a localized signal-processing
pass rather than a global neural decoder. Haptics is best-served
across the contact-mediated portion of the cycle: contact
establishment in phase 2, in-contact transport in phase 3, and
release in phase 4. The modality is moderately coupled to active
perception: a non-empty $D_{h}(\jointvec)$ requires that phase 2
has closed the contact and the fingertip is touching the object,
so haptics depends on the controller's ability to bring the
fingertip into contact, but once contact is established the
sensor reading is determined by the local geometry rather than by
continued reorientation of the arm.

\paragraph{Modality complementarity.}
The three modalities are complementary across the five-phase
cycle. Vision dominates the out-of-contact reaches (phases 1 and
5), where the relevant scene quantities lie outside the
fingertip's contact patch and the latency budget is large enough
to absorb the moderate $\obsmap_{v}$ pipeline. Proprioception
dominates the bang arcs (phases 2 and 4) at the
$\swsurf$-crossings, where the timeliness condition
\eqref{eq:timeliness} is tightest and only the encoder-bus rate
is fast enough to close the loop on the $\rcontact$-threshold
detection. Haptics dominates the contact-mediated portion of the
in-contact transport (phase 3), where the contact patch carries
exactly the local-equilibrium information that the bi-level
inner program of \cref{def:bilevel} requires. The Stage IV
bi-manual block-place + nail-hold + hammer task of
\cref{sec:experiments} exercises all three modalities: vision
locates the block, the nail, and the hammer for the pre-grasp and
retreat reaches; proprioception detects each $\swsurf$-crossing
during contact establishment and release with the block, the
nail, and the hammer; and haptics carries the contact-patch
information during the in-contact transport segments of each
sub-task. The three observation maps
$\obsmap_{v}, \obsmap_{p}, \obsmap_{h}$ together specialize the
single-sensor abstraction of \cref{def:obsmap} to the
multi-sensor architecture that the controller of
\cref{sec:plan-training} will consume.
